\section{A política}
\label{sec:politica}

\subsection{A lei e a sua cronologia}

A Lei Estadual nº~16.820, de 8 de janeiro de 2019, acrescentou o art.~28-B à
lei cearense de agrotóxicos (Lei nº~12.228/1993) e vedou a pulverização aérea de
agrotóxicos na agricultura em todo o estado, com multa ao infrator. Entrou em
vigor na publicação, em 9 de janeiro. A Tabela~\ref{tab:cronologia} reúne os
marcos.

\begin{table}[ht]
\centering\small
\caption{Cronologia da política}
\label{tab:cronologia}
\begin{tabularx}{\textwidth}{@{}lL@{}}
\toprule
Data & Fato \\
\midrule
20/11/2009 & Limoeiro do Norte proíbe a pulverização aérea por lei municipal (Lei 1.478). \\
21/04/2010 & Assassinato de José Maria Filho, o Zé Maria do Tomé. \\
20--26/05/2010 & A Câmara de Limoeiro revoga a proibição, por voto de minerva, dentro da lei da política ambiental do município. \\
24/02/2015 & Apresentado na Assembleia Legislativa o Projeto de Lei 18/2015. \\
18/12/2018 & O projeto é aprovado por unanimidade. \\
09/01/2019 & A Lei 16.820 entra em vigor. \\
05/2023 & O Supremo Tribunal Federal julga a lei constitucional (ADI 6137, unânime). \\
19/12/2024 & A Lei 19.135 abre exceção para a aplicação por drones. \\
\bottomrule
\end{tabularx}
\end{table}

Três traços da cronologia importam para a avaliação. A lei é estadual e
simultânea: não há adoção escalonada entre municípios, e a variação que
identifica o efeito vem da fronteira com um estado que não proibiu. O projeto
tramitou por quase quatro anos, de modo que os produtores tiveram notícia da
proibição muito antes de ela valer. E a exceção de dezembro de 2024 encerra o
período em que o tratamento é a proibição total; a janela de análise termina
nela.

O precedente municipal de Limoeiro do Norte durou seis meses e, segundo o
registro do processo criminal e o trabalho de campo da época, não interrompeu
os voos \citep[p.~536--542]{teixeira2011rebulico}. Entre 2010 e 2018 não houve
proibição vigente na Chapada.

\subsection{A Chapada do Apodi e a prática que a lei proibiu}

A Chapada do Apodi é um planalto calcário na divisa entre o Ceará e o Rio
Grande do Norte. Do lado cearense, o Perímetro Irrigado Jaguaribe--Apodi, nos
municípios de Limoeiro do Norte e Quixeré, concentrou a partir dos anos 2000 a
fruticultura de exportação, com empresas nacionais e transnacionais e pequenos
produtores a elas associados. Do lado potiguar, a mesma fruticultura irrigada
ocupa Baraúna, Apodi e o entorno de Mossoró. Os dois lados partilham o clima
semiárido, a quadra chuvosa de fevereiro a maio e o Aquífero Jandaíra.

O avião servia à banana. O alvo era a sigatoka-amarela, doença fúngica que
avança com a umidade, e por isso as aplicações se concentravam na quadra
chuvosa. A calda era de cerca de 25 litros por hectare: meio litro de
fungicida, cinco de óleo e vinte de água. Os produtos eram triazóis e
estrobilurinas, entre eles difenoconazol, propiconazol, tebuconazol e
epoxiconazol \citep[p.~526]{teixeira2011rebulico}. O químico, portanto, não é o
glifosato que domina a literatura econômica sobre agrotóxicos.

Sobre a escala, as fontes divergem, e a divergência é informativa. O trabalho de
campo de 2009 observou um contrato de 1.000 hectares, dos quais 820 de duas
empresas e 180 de vinte pequenos produtores parceiros, e uma primeira aplicação
sobre 637 hectares \citep[p.~171--173]{marinho2011dimensao}. Uma nota técnica do
setor produtivo, de 2016, fala em cerca de 1.000 hectares pulverizados entre
2003 e 2015, ou 2,27\% da banana do estado \citep{tramas2016parecer}. O número
que circula no debate público, de quatro milhões de litros de calda em dez
anos, é uma conta sobre toda a área de banana da Chapada e deve ser lido como
teto \citep[p.~527]{teixeira2011rebulico}.

O Censo Agropecuário de 2006 dá a linha de base por município
(Tabela~\ref{tab:censo}). O avião era usado em Limoeiro do Norte e Quixeré,
sobretudo por estabelecimentos não familiares, e não aparece no lado potiguar
da Chapada.

\begin{table}[ht]
\centering\small
\caption{Estabelecimentos que usaram agrotóxico, por equipamento de aplicação,
2006}
\label{tab:censo}
\begin{tabular}{@{}lrrrrr@{}}
\toprule
 & \multicolumn{3}{c}{Por aeronave} & Tração & \\
\cmidrule(lr){2-4}
Município & Total & Não familiar & Familiar & mecânica & Costal \\
\midrule
Limoeiro do Norte (CE) & 18 & 14 & 4 & 56 & 871 \\
Quixeré (CE)           &  9 &  8 & 1 & 15 & 506 \\
Russas (CE)            &  0 &  0 & 0 & 18 & 2.020 \\
Apodi (RN)             &  0 &  0 & 0 & 30 & 1.503 \\
Baraúna (RN)           &  0 &  0 & 0 & 107 & 418 \\
Mossoró (RN)           &  0 &  0 & 0 & 15 & 546 \\
\midrule
Ceará, total           & 36 & 24 & 12 & & \\
Rio Grande do Norte, total & 5 & 0 & 5 & & \\
\bottomrule
\multicolumn{6}{@{}p{0.95\textwidth}@{}}{\footnotesize Fonte: IBGE, Censo
Agropecuário 2006, tabela SIDRA 2069, extraída em 16/11/2026. Um
estabelecimento pode declarar mais de um equipamento. O Censo de 2017 não
repetiu a pergunta.}
\end{tabular}
\end{table}

\subsection{O que não se sabe sobre os anos anteriores à lei}
\label{sec:quando}

A pulverização aérea na Chapada está documentada com medida até 2010. Para os
anos seguintes há relatos, e eles sugerem que a prática rareou ou cessou com a
seca de 2012 a 2017. Uma moradora disse em 2016 que desde 2012 não havia novas
aplicações, por falta de chuva; a coordenadora do principal estudo
epidemiológico da região afirmou no mesmo ano que a pulverização aérea não
vinha acontecendo em função da seca; e uma entrevistada de \citet[p.~115]{aguiar2017}
referiu-se a ela, em 2017, como coisa de alguns anos antes. Em sentido
contrário, o setor produtivo litigou contra a lei até o Supremo, e um produtor
da Chapada descreveu à imprensa, em 2021, a troca do avião pela aplicação
manual.

Essa incerteza não é um defeito do projeto; é parte do que ele se propõe a
resolver. A Instrução Normativa nº~2/2008 do Ministério da Agricultura obriga
todo operador aeroagrícola a entregar relatório mensal com município, cultura e
área aplicada. O Decreto nº~4.074/2002 exige que o receituário agronômico
registre a modalidade de aplicação. Esses registros medem o que a lei removeu, e
o projeto os solicita pela Lei de Acesso à Informação. A Seção~\ref{sec:cenarios}
diz o que o desenho estima em cada resposta possível.

\subsection{Fiscalização e o regime no estado vizinho}

A norma federal já impunha distâncias mínimas à aplicação aérea: 500 metros de
povoações e 250 metros de mananciais. A literatura da Chapada descreve o
descumprimento dessas distâncias \citep{teixeira2011rebulico,
tramas2016parecer}. A proibição total troca uma regra difícil de fiscalizar por
uma simples: avião aplicando é infração.

O Rio Grande do Norte permaneceu sob o regime federal. Não há lei estadual
potiguar que proíba a pulverização aérea; um projeto nesse sentido foi
apresentado na Assembleia em novembro de 2019 e não se converteu em lei. No
levantamento legislativo de \citet{olivete2025}, o Ceará segue como o único
estado com proibição total; os demais têm normas de registro e fiscalização ou
leis municipais.
